% =============================================================================
% Supplementary Section S2 : Model comparison
% Response to Reviewer 2 major point 3
% =============================================================================

\section*{S2. Formal model comparison: is the quadratic form required?}
\label{sec:S2}

\subsection*{Motivation}

A legitimate methodological question is whether the
quadratic model used in the main analysis may impose a particular
functional form rather than being required by the data. This motivates
a formal model comparison using information criteria (AIC,
BIC) to assess whether other functional forms (flat, linear, broken-line,
spline) might describe the data equally well or better.

We address this concern by fitting five competing functional forms to
the SPARC sample ($N = 181$ galaxies) and comparing them using
information criteria with appropriate penalisation for model complexity.

\subsection*{Models tested}

The five candidate models are:

\begin{itemize}
    \item \textbf{Flat} ($k = 1$): $\Delta(\rho) = c$ (constant, null hypothesis)
    \item \textbf{Linear} ($k = 2$): $\Delta(\rho) = b\rho + c$
    \item \textbf{Quadratic} ($k = 3$): $\Delta(\rho) = a\rho^2 + b\rho + c$ (used in main text)
    \item \textbf{Broken-line} ($k = 4$): two linear segments joined at a free breakpoint
    \item \textbf{Cubic spline} ($k = 5$): natural cubic spline with two internal knots
\end{itemize}

where $\rho$ is the environmental density proxy, $k$ is the number of
free parameters, and $c$, $b$, $a$, breakpoint locations and spline
knots are fitted parameters.

For each model we compute the maximum likelihood, the Akaike Information
Criterion (AIC), the Bayesian Information Criterion (BIC), Akaike
weights $w_i$ (relative probability of each model being the best one in
the candidate set), and the coefficient of determination $R^2$.

\subsection*{Results}

\begin{table}[htbp]
\centering
\caption{Model comparison on the SPARC sample ($N = 181$). Lower
$\Delta\mathrm{AIC}$ and $\Delta\mathrm{BIC}$ indicate better fit (after
penalisation for model complexity). All non-monotonic models (quadratic,
broken-line, spline) are strongly preferred over monotonic alternatives
(flat, linear).}
\label{tab:S2_model_comparison}
\begin{tabular}{lrrrrrr}
\toprule
\textbf{Model} & $k$ & $\Delta\mathrm{AIC}$ & $\Delta\mathrm{BIC}$ & $w_i$ (AIC) & $R^2$ \\
\midrule
Cubic spline & 5 & 0.00 & 2.36 & 0.602 & 0.287 \\
Broken-line & 4 & 0.84 & 0.00 & 0.396 & 0.276 \\
Quadratic & 3 & 11.14 & 7.11 & 0.002 & 0.225 \\
Linear & 2 & 39.95 & 32.71 & $1.3 \times 10^{-9}$ & 0.081 \\
Flat & 1 & 53.24 & 42.81 & $1.7 \times 10^{-12}$ & 0.000 \\
\bottomrule
\end{tabular}
\end{table}

\begin{figure}[htbp]
\centering
\includegraphics[width=\textwidth]{figureS2_model_comparison.png}
\caption{\textbf{Model comparison on the SPARC sample.} Panel A shows
$\Delta\mathrm{AIC}$ values for the five candidate models, with the
red dashed line indicating the $\Delta\mathrm{AIC} = 10$ threshold for
strong evidence against a model. Panel B shows the same for
$\Delta\mathrm{BIC}$. Panel C shows the Akaike weights $w_i$ on a
logarithmic scale.}
\label{fig:S2}
\end{figure}

\subsection*{Interpretation}

The results lead to two complementary conclusions:

\paragraph{(i) Non-monotonic dependence is strongly favoured by the data within the tested model classes.}
All three non-monotonic models (quadratic, broken-line, spline) are
strongly preferred over the monotonic alternatives (linear and flat).
Specifically, the simplest non-monotonic model (quadratic) outperforms
the linear model by $\Delta\mathrm{AIC} = 28.8$ and $\Delta\mathrm{BIC} = 25.6$,
and outperforms the flat model by $\Delta\mathrm{AIC} = 42.1$ and
$\Delta\mathrm{BIC} = 35.7$. These differences correspond to relative
likelihoods of order $10^{-7}$ to $10^{-10}$ against the monotonic
hypotheses. The non-monotonic dependence is therefore strongly favoured within the tested model classes
by the data and therefore does not depend on the specific choice of a quadratic form.

\paragraph{(ii) The quadratic form is slightly outperformed by more flexible models.}
The cubic spline and broken-line models, which have more free parameters,
provide slightly better fits than the pure quadratic ($\Delta\mathrm{AIC} = 11$,
$\Delta\mathrm{BIC} = 7$). This suggests that the true underlying
functional form may be more complex than a pure parabola, e.g., with
asymmetric arms or a sharper transition between regimes. However, the
difference between quadratic and the more flexible models is much smaller
than the difference between quadratic and the monotonic alternatives.

\subsection*{Why we retain the quadratic model in the main text}

We retain the quadratic parameterization in the main text for three
reasons:

\begin{enumerate}
    \item \textbf{Physical interpretability}: the single curvature
    coefficient $a$ has a direct interpretation as the strength of the
    antagonistic Q-R coupling in the QO+R framework. The broken-line
    and spline parameters do not carry comparable physical meaning.

    \item \textbf{Parsimony}: among models that capture the qualitative
    non-monotonic pattern, the quadratic has the fewest free parameters
    ($k = 3$) and the simplest interpretation.

    \item \textbf{Comparability with other surveys}: a single curvature
    coefficient $a$ allows direct comparison across datasets with
    different sample sizes (SPARC, ALFALFA, Little THINGS), which would
    be impractical with multi-parameter spline or breakpoint models.
\end{enumerate}

We acknowledge in the main text and in this supplementary material that
more flexible descriptive models fit slightly better, and we encourage
future work to investigate the precise functional form of the
environmental dependence as larger samples become available.

\subsection*{Conclusion}

The non-monotonic dependence of BTFR residuals on environment is
\textbf{strongly favoured within the tested model classes} by the SPARC data, with monotonic alternatives
ruled out at $\Delta\mathrm{AIC} > 39$ and $\Delta\mathrm{BIC} > 32$.
The quadratic form is retained in the main text for its physical
interpretability and parsimony, while more flexible models (spline,
broken-line) provide marginally better statistical fits that may guide
future theoretical refinements.
